% Source-aligned English for printed p.190.
\clearpage
\SourcePage{190}

\LayerTarget{tohoku:corollary:4.2.4:1}
\LayerTarget{tohoku:eq:4.2.10}
\[
\tag{4.2.10}
\mathbf{Ext}_O^n(A,B)=H^n(\mathbf{Hom}_O(\mathbf{L},B)),
\]
and hence

\LayerTarget{tohoku:eq:4.2.10-bis}
\[
\tag{4.2.10 bis}
\mathbf{Ext}_O^n(A,B)(x)=
\mathrm{Ext}_{O(x)}^n(A(x),B(x)).
\]

\LayerTarget{tohoku:proof:corollary-4.2.4:1}
Indeed, it suffices to apply Proposition 2.5.1.

Applying Proposition 2.5.4 now gives:

\LayerTarget{tohoku:corollary:4.2.4:2}
\noindent\textsc{Corollary 2.} \textit{Under the preceding hypotheses, the
spectral functor of Theorem 4.2.1 is isomorphic to the spectral functor
$\Pi\Gamma(\mathbf{Hom}_O(\mathbf{L},B))$. In particular,}
\[
\mathrm{Ext}_O^n(X;A,B)=
\mathbb{H}^n\Gamma(\mathbf{Hom}_O(\mathbf{L},B)),
\]
\textit{where the right-hand side is the n-th hyperhomology group of
$\Gamma$ with respect to the complex $\mathbf{Hom}_O(\mathbf{L},B)$
(cf. 2.4).}

\LayerTarget{tohoku:construction:canonical-resolvent-spectral-sequence}
We also have a resolvent functor for $\Gamma$, for example the functor
$\Gamma C(F)$, where $C(F)$ is the canonical resolution of $F$ (cf. 2.3).
Thus the spectral sequence $\Pi\Gamma(\mathbf{Hom}_O(\mathbf{L},B))$ becomes
explicit: it is the first spectral sequence of the bicomplex
$\Gamma C(\mathbf{Hom}_O(\mathbf{L},B))$, with the degree coming from $C$
taken as the first degree. There is a functorial isomorphism

\LayerTarget{tohoku:eq:4.2.11}
\[
\tag{4.2.11}
C(\mathbf{Hom}_O(\mathbf{L},B))=
\mathbf{Hom}_O(\mathbf{L},C(B)).
\]

Indeed, for every $O$-module $L$ there is a functorial homomorphism

\LayerTarget{tohoku:eq:4.2.12}
\[
\tag{4.2.12}
C(\mathbf{Hom}_O(L,B))\longrightarrow\mathbf{Hom}_O(L,C(B)).
\]

\LayerTarget{tohoku:proof:equations-4.2.11-4.2.12}
Here, on the right, $C(B)$ is regarded in the evident way as an $O$-module.
The map is defined inductively on the dimension of the components of $C$ from
the natural homomorphism
$C^0(\mathbf{Hom}_O(L,B))\longrightarrow\mathbf{Hom}_O(L,C^0(B))$, itself
deduced from the natural homomorphisms
$\mathbf{Hom}_O(L,B)(x)\longrightarrow
\mathrm{Hom}_{O(x)}(L(x),B(x))$. The latter is an isomorphism when $L$ is
pseudo-coherent (Proposition 4.1.1), so in that case (4.2.12) is an
isomorphism; this proves (4.2.11). It follows that
$\Gamma C(\mathbf{Hom}_O(\mathbf{L},B))=
\mathbf{Hom}_O(\mathbf{L},C(B))$, whence:

\LayerTarget{tohoku:corollary:4.2.4:3}
\noindent\textsc{Corollary 3.} \textit{Let $\mathbf{L}$ be as in Proposition
4.2.4. The spectral functor of Theorem 4.2.1 is given by the first spectral
sequence $\Gamma\mathbf{Hom}_O(\mathbf{L},C(B))$ of the bicomplex
$\mathbf{Hom}_O(\mathbf{L},C(B))$, where $C(B)$ is the canonical resolution
(cf. 2.3) of the abelian sheaf $B$, whose grading is taken as the first degree
of the bicomplex. In particular,}
\[
\mathrm{Ext}_O^n(X;A,B)=
H^n(\mathbf{Hom}_O(\mathbf{L},C(B))).
\]

\LayerTarget{tohoku:remark:ext-with-supports}
\noindent\textsc{Remarks.} 1) One may also define the functors
$\mathrm{Ext}_{O,\varphi}^p(X;A,B)$ as the right derived functors, with
respect to $B$, of
$\mathrm{Hom}_{O,\varphi}(A,B)=
\Gamma_\varphi(\mathbf{Hom}_O(A,B))$, where $\varphi$ is an antifilter of
closed subsets of $X$. The preceding developments remain valid. We shall not
need them.

\LayerTarget{tohoku:remark:finiteness-in-theorem-4.2.2}
2) It is essential in Theorem 4.2.2 that $A$ satisfy a finiteness condition.
For example, if $A=O^{(I)}$ with $I$ infinite, then the right-hand side of
(4.2.8) vanishes for every $x$ and every $p>0$, but in general one does
\textit{not} have $\mathbf{Ext}_O^1(A,B)=0$, since the functor
$B\longmapsto\mathbf{Hom}_O(O^{(I)},B)=B^I$ is not exact in general.

\TohokuContinuation{191}
